% Generated draft: requires mathematical and source review.
\documentclass[11pt]{article}
\usepackage[margin=25mm]{geometry}
\usepackage{fontspec}
\setmainfont{Latin Modern Roman}
\usepackage{amsmath,amssymb,mathtools,unicode-math}
\setmathfont{Latin Modern Math}
\usepackage{enumitem,needspace,xurl,hyperref}
\setlength{\parindent}{0pt}
\setlength{\parskip}{5pt}
\setlength{\emergencystretch}{3em}
\newcommand{\N}{\mathbb{N}}
\newcommand{\Z}{\mathbb{Z}}
\newcommand{\Q}{\mathbb{Q}}
\newcommand{\R}{\mathbb{R}}
\newcommand{\C}{\mathbb{C}}
\newcommand{\F}{\mathbb{F}}
\newcommand{\A}{\mathbb{A}}
\newcommand{\PP}{\mathbb{P}}
\newcommand{\E}{\mathbb{E}}
\newcommand{\eps}{\varepsilon}
\newcommand{\dd}{\,\mathrm{d}}
\begin{document}
\section*{UnsolvedMath ID: 701706}
\textbf{Source problem code:} catalogue entry
\subsection{Definitions and mathematical statement}
Let \(p\) be a prime and let \(P\) be a nontrivial finite \(p\)-group. For this draft, a nontrivial partition of \(P\) means a family \(\mathcal P\) of at least two nonidentity proper subgroups such that
\[
P=\bigcup_{U\in\mathcal P}U,
\qquad U\cap V=\{1\}\quad\text{for distinct }U,V\in\mathcal P.
\]
Define the Hughes subgroup by
\[
H_p(P)=\left\langle g\in P\mid g^p\ne 1\right\rangle.
\]
An automorphism \(\alpha\in\operatorname{Aut}(P)\) is fixed-point-free if
\[
\{x\in P\mid \alpha(x)=x\}=\{1\}.
\]
The statement that \(P\) has exponent \(p\) means that every element satisfies \(x^p=1\) and some element has order \(p\).

The question asks whether, for every prime \(p\) and every nontrivial finite \(p\)-group \(P\) admitting a nontrivial partition, the following implication holds:
\[
\left(\exists\text{ a prime }r\;\exists\alpha\in\operatorname{Aut}(P):
\operatorname{ord}(\alpha)=r
\ \text{and}\ \{x\in P\mid\alpha(x)=x\}=\{1\}\right)
\Longrightarrow \operatorname{exp}(P)=p.
\]
No relation between \(r\) and \(p\) is imposed in the supplied statement.

The supplied record identifies the partition hypothesis with \(H_p(P)\ne P\). Its alternate formulation therefore replaces the partition hypothesis above by this proper-Hughes-subgroup condition. The source does not specify its convention for a nontrivial partition. Under the convention used here, the cyclic group of order \(p\) has proper Hughes subgroup but has no nontrivial partition. Thus the asserted equivalence requires a boundary qualification or a different partition convention; the source's intended convention remains unresolved.
\subsection{Short English statement}
Must a finite \(p\)-group with a nontrivial partition have exponent \(p\) whenever it admits a fixed-point-free automorphism of prime order? The source also gives the hypothesis as \(H_p(P)\ne P\); the precise equivalence of these hypotheses has an unresolved small-order convention.
\subsection{Sources}
\begin{itemize}
\item Local UnsolvedMath registry, permanent ID \(701706\), exact source problem code \texttt{catalogue entry}, title \emph{Kourovka Notebook Problem 20.48}. Supplied upstream revision: \texttt{d371ee917ac9647fbb06cd6fc529deec6f3dda82}.
\item Attributed in the supplied record to M. Lewis, proposed in 2022; original problem number 20.48 in \emph{The Kourovka Notebook (Unsolved Problems in Group Theory)}, 20th issue (2022), with the 21st issue edition (2026), \texttt{arXiv:1401.0300v48}, also identified as a source. The supplied collection description names E. I. Khukhro and V. D. Mazurov as editors.
\item The registry labels this problem \emph{open} and gives an update timestamp of \texttt{2026-10-16T14:00:00Z}. These are supplied registry metadata, not an independently verified assessment of the literature. The referenced editions and arXiv version have not been independently consulted for this draft.
\end{itemize}
\end{document}
